\documentclass[11pt,letterpaper]{article}
\usepackage[lmargin=1in,rmargin=1in,tmargin=1in,bmargin=1in]{geometry}
\usepackage{wrapfig}
\usepackage{colortbl}
% -------------------
% Packages
% -------------------
\usepackage{
	amsmath,			% Math Environments
	amssymb,			% Extended Symbols
	enumerate,		    % Enumerate Environments
	graphicx,			% Include Images
	lastpage,			% Reference Lastpage
	multicol,			% Use Multi-columns
	multirow,			% Use Multi-rows
    ulem,                % Make underlines
    subcaption,
    longtable,
    booktabs,
    float,
    array,
    calc,
    charter,
    hyperref,
    pgfplots,
    xcolor
}

% -------------------
% Font
% -------------------
\usepackage[en-US]{datetime2}
\newlength{\mycellsize} % define a custom length for cell size
\setlength{\mycellsize}{1.5cm} % set the cell size to 2cm
\usepackage[T1]{fontenc}
\pgfplotsset{compat=1.18}
% -------------------
% Heading Commands
% -------------------
\newcommand{\class}{MA205: Multivariable Calculus}
\newcommand{\term}{AY27-1}
\newcommand{\instructor}{MAJ John Borger (TH255)}
\newcommand{\head}[1]{%
\thispagestyle{empty} % Ensures no header or footer
\vspace*{-0.5in}
\noindent
\begin{tabular}{@{} m{1.5cm}@{} m{3cm}@{} c@{}}
\includegraphics[height=1.5cm]{Images/Mule.png} & &

\begin{tabular}{@{}l@{}}
  \textbf{Name:} \makebox[10.5cm]{\hrulefill} \\
  \textbf{#1} %\textbf{August 20, 2024}
  \\
  \textbf{\class} \\
  \textbf{Instructor: \instructor, \term} \\
\end{tabular}
\end{tabular} \\
\rule[2ex]{\textwidth}{2pt} %
\vspace*{-0.8cm}
}

% -------------------
% Commands
% -------------------
\newcommand{\prob}{\noindent\textbf{Problem. }}
\newcounter{problem}
\newcommand{\problem}{
	\stepcounter{problem}%
	\noindent \textbf{Problem \theproblem. }%
}
\newcommand{\pointproblem}[1]{
	\stepcounter{problem}%
	\noindent \textbf{Problem \theproblem.} (#1 points)\,%
}
\newcommand{\pspace}{\par\vspace{\baselineskip}}
\newcommand{\ds}{\displaystyle}


% -------------------
% Header & Footer
% -------------------
\usepackage{fancyhdr}

\fancypagestyle{pages}{
	%Headers
	\fancyhead[L]{}
	\fancyhead[C]{}
	\fancyhead[R]{}
\renewcommand{\headrulewidth}{0pt}
	%Footers
	\fancyfoot[L]{}
	\fancyfoot[C]{}
	\fancyfoot[R]{\thepage}
\renewcommand{\footrulewidth}{0.0pt}
}
\headheight=0pt
\footskip=14pt

\pagestyle{pages}

% Let's me turn on and off showing solutions
% Instructor copy: show solutions in red. For a student handout, comment the next line
% and uncomment \showsolutionsfalse.
\newif\ifshowsolutions
\showsolutionstrue
%\showsolutionsfalse

% Solution macro: red text only when \ifshowsolutions is true
\newcommand{\sol}[1]{%
\ifshowsolutions
\par\vspace{0.3em}{\color{red}#1}\par
\fi
}


% -------------------
% Content
% -------------------
\begin{document}
 
\head{Lesson 32 \& 33 Applications of Double Integrals}

\noindent\textbf{Learning Objectives:}
\vspace{-2mm}
\begin{enumerate}
\setlength{\itemsep}{0pt} % Adjusts the space between items
\setlength{\parskip}{0pt} % Adjusts the space between paragraphs

    \item Determine the mass and center of mass of a lamina of variable density using double integrals.
    \item Determine the probability of an event and the expected values of a joint density function using double integrals.

\end{enumerate}



\hrule \vspace{0.1cm}
\noindent \textbf{Reading:} Supplemental Reading : \S 15.4
\vspace{0.1cm}
\hrule

\section{Warm Up}


    \subsection{}
    
    Evaluate the integral $\iint_D e^{-x^2-y^2} \, dA$ where $D$ is the region bounded by the semicircle $x=+\sqrt{4-y^2}$ and $x=0.$

\sol{%
\textbf{Solution.}
The curve $x=\sqrt{4-y^2}$ is the right half of $x^2+y^2=4$, so $D$ is the right semicircle of radius $2$ centered at the origin.\\[0.4em]
Switch to polar: $x=r\cos\theta$, $y=r\sin\theta$, $dA=r\,dr\,d\theta$.\\
Then $e^{-x^2-y^2}=e^{-r^2}$, with $0\leq r\leq 2$ and $-\frac{\pi}{2}\leq\theta\leq\frac{\pi}{2}$.
\begin{align*}
\iint_D e^{-x^2-y^2}\,dA
&= \int_{-\pi/2}^{\pi/2}\int_0^{2} e^{-r^2}\,r\,dr\,d\theta.
\end{align*}
Inner integral: let $u=r^2$, so $du=2r\,dr$.
\begin{align*}
\int_0^{2} e^{-r^2}r\,dr
&= \frac{1}{2}\int_0^{4} e^{-u}\,du
= \frac{1}{2}\bigl[-e^{-u}\bigr]_0^{4}
= \frac{1}{2}\bigl(1-e^{-4}\bigr).
\end{align*}
Outer integral:
\begin{align*}
\int_{-\pi/2}^{\pi/2}\frac{1}{2}\bigl(1-e^{-4}\bigr)\,d\theta
&= \frac{1}{2}\bigl(1-e^{-4}\bigr)\cdot\pi
= \frac{\pi}{2}\bigl(1-e^{-4}\bigr).
\end{align*}
}

\vfill 

\pagebreak


\subsection{}

Convert to polar coordinates and solve.
        $\ds \int_0^{1/2}\int_{\sqrt{3}y}^{\sqrt{1-y^2}} xy^2 \, dx\, dy$.

\sol{%
\textbf{Solution.}
Sketch the region: the left boundary $x=\sqrt{3}y$ is the ray $\theta=\pi/6$, and the right boundary $x=\sqrt{1-y^2}$ is the unit circle $r=1$. With $y$ from $0$ to $1/2$, we have $0\leq\theta\leq\pi/6$ and $0\leq r\leq 1$.\\[0.4em]
In polar: $x=r\cos\theta$, $y=r\sin\theta$, $dA=r\,dr\,d\theta$, so
\begin{align*}
xy^2\,dA &= (r\cos\theta)(r\sin\theta)^2\cdot r\,dr\,d\theta = r^4\cos\theta\sin^2\theta\,dr\,d\theta.
\end{align*}
\begin{align*}
\int_0^{\pi/6}\int_0^{1} r^4\cos\theta\sin^2\theta\,dr\,d\theta
&= \int_0^{\pi/6}\cos\theta\sin^2\theta\Bigl[\frac{r^5}{5}\Bigr]_0^{1}\,d\theta
= \frac{1}{5}\int_0^{\pi/6}\cos\theta\sin^2\theta\,d\theta.
\end{align*}
Let $u=\sin\theta$, $du=\cos\theta\,d\theta$. When $\theta=0$, $u=0$; when $\theta=\pi/6$, $u=1/2$.
\begin{align*}
\frac{1}{5}\int_0^{1/2} u^2\,du
&= \frac{1}{5}\cdot\frac{u^3}{3}\Big|_0^{1/2}
= \frac{1}{15}\cdot\frac{1}{8}
= \frac{1}{120}.
\end{align*}
}

\vfill 


\subsection{}

Find the area enclosed by one loop of the four-leaved rose $r=\cos(2\theta)$ shown below.\\

    \includegraphics[scale=1]{Images/L34 4LeafClover.png}

\sol{%
\textbf{Solution.}
One loop of $r=\cos(2\theta)$ occurs where $\cos(2\theta)\geq 0$ and $r$ returns to $0$, namely $-\pi/4\leq\theta\leq\pi/4$.
\begin{align*}
A &= \frac{1}{2}\int_{-\pi/4}^{\pi/4} r^2\,d\theta
= \frac{1}{2}\int_{-\pi/4}^{\pi/4}\cos^2(2\theta)\,d\theta
= \frac{1}{2}\int_{-\pi/4}^{\pi/4}\frac{1+\cos(4\theta)}{2}\,d\theta
= \frac{1}{4}\int_{-\pi/4}^{\pi/4}\bigl(1+\cos(4\theta)\bigr)\,d\theta.
\end{align*}
\begin{align*}
A &= \frac{1}{4}\Bigl[\theta+\frac{1}{4}\sin(4\theta)\Bigr]_{-\pi/4}^{\pi/4}
= \frac{1}{4}\Bigl(\frac{\pi}{4}-\Bigl(-\frac{\pi}{4}\Bigr)\Bigr)
= \frac{1}{4}\cdot\frac{\pi}{2}
= \frac{\pi}{8}.
\end{align*}
}

\vfill 

\pagebreak 


 \section{Probability}
Suppose two numbers $x$ and $y$ are randomly generated by a computer such that ${0\leq x\leq 1}$ and $0\leq y\leq 2$. The joint density function is given below.
$$ f(x,y) = \begin{cases}
Cx(1+y)\qquad \text{ if $ 0\leq x\leq 1$ and $ 0\leq y\leq 2$}\\
0 \qquad\qquad\qquad\qquad \text{ otherwise} \end{cases} $$
 
 
 \begin{enumerate}[a)]
\item Find the value of $C$ such that $f(x,y)$ is a joint density function.
\sol{%
\textbf{Solution.}
A joint density must satisfy $\iint_{\mathbb{R}^2} f(x,y)\,dA=1$:
\begin{align*}
\int_0^1\int_0^2 Cx(1+y)\,dy\,dx
&= C\int_0^1 x\,dx\int_0^2(1+y)\,dy
= C\cdot\frac{1}{2}\cdot\Bigl[y+\frac{y^2}{2}\Bigr]_0^2
= C\cdot\frac{1}{2}\cdot 4
= 2C.
\end{align*}
Set $2C=1$, so $\displaystyle C=\frac{1}{2}$.
}
\vfill
\item Using your $C$ value from part $a)$, find $P( .5 \leq Y\leq 1.5)$.
\sol{%
\textbf{Solution.}
With $C=1/2$,
\begin{align*}
P(0.5\leq Y\leq 1.5)
&= \int_0^1\int_{0.5}^{1.5}\frac{1}{2}x(1+y)\,dy\,dx
= \frac{1}{2}\int_0^1 x\,dx\int_{0.5}^{1.5}(1+y)\,dy.
\end{align*}
\begin{align*}
\int_{0.5}^{1.5}(1+y)\,dy
&= \Bigl[y+\frac{y^2}{2}\Bigr]_{1/2}^{3/2}
= \Bigl(\frac{3}{2}+\frac{9}{8}\Bigr)-\Bigl(\frac{1}{2}+\frac{1}{8}\Bigr)
= \frac{21}{8}-\frac{5}{8}
= 2.
\end{align*}
\begin{align*}
P &= \frac{1}{2}\cdot\frac{1}{2}\cdot 2 = \frac{1}{2}.
\end{align*}
}
\vfill
\item What is the probability that the sum of $x$ and $y$ is less than 1?
\sol{%
\textbf{Solution.}
We want $P(X+Y<1)$ over the rectangle $[0,1]\times[0,2]$. That region is $0\leq x\leq 1$, $0\leq y\leq 1-x$ (since $1-x\leq 1<2$).
\begin{align*}
P(X+Y<1)
&= \int_0^1\int_0^{1-x}\frac{1}{2}x(1+y)\,dy\,dx
= \frac{1}{2}\int_0^1 x\Bigl[y+\frac{y^2}{2}\Bigr]_0^{1-x}\,dx \\
&= \frac{1}{2}\int_0^1 x\Bigl((1-x)+\frac{(1-x)^2}{2}\Bigr)\,dx
= \frac{1}{2}\int_0^1 x\cdot\frac{2(1-x)+(1-x)^2}{2}\,dx \\
&= \frac{1}{4}\int_0^1 x\bigl(2-2x+1-2x+x^2\bigr)\,dx
= \frac{1}{4}\int_0^1\bigl(3x-4x^2+x^3\bigr)\,dx \\
&= \frac{1}{4}\Bigl[\frac{3}{2}x^2-\frac{4}{3}x^3+\frac{1}{4}x^4\Bigr]_0^1
= \frac{1}{4}\Bigl(\frac{3}{2}-\frac{4}{3}+\frac{1}{4}\Bigr)
= \frac{1}{4}\cdot\frac{5}{12}
= \frac{5}{48}.
\end{align*}
}
\end{enumerate}

\vfill


\newpage

\section{Density and Center of Mass}
\subsection{Density ``Understand''}

Let $\rho $ be density, $V$ be volume, and $m$ be mass.
\begin{align*}
    \rho= \frac{m}{V} && m=\rho V
\end{align*}

\vfill 
\sol{%
For a planar lamina with variable density $\rho(x,y)$,
\begin{align*}
    m=\iint_D \rho(x,y) \, dA.
\end{align*}
}


\subsection{Example ``Do''}
Find the mass of a triangular lamina with vertices $(0,0),(1,0)$ and $(0,2)$ and density function $\rho(x,y)=1+3x+y$.

\sol{%
\textbf{Solution.}
The hypotenuse from $(1,0)$ to $(0,2)$ is $y=2-2x$. So $0\leq x\leq 1$, $0\leq y\leq 2-2x$.
\begin{align*}
m &= \int_0^1\int_0^{2-2x}(1+3x+y)\,dy\,dx
= \int_0^1\Bigl[y+3xy+\frac{y^2}{2}\Bigr]_0^{2-2x}\,dx.
\end{align*}
Let $h=2-2x$. Then
\begin{align*}
y+3xy+\frac{y^2}{2}\Big|_0^{h}
&= h+3xh+\frac{h^2}{2}
= (2-2x)+3x(2-2x)+\frac{(2-2x)^2}{2} \\
&= 2-2x+6x-6x^2+2-4x+2x^2
= 4-4x^2.
\end{align*}
\begin{align*}
m &= \int_0^1(4-4x^2)\,dx
= \Bigl[4x-\frac{4x^3}{3}\Bigr]_0^1
= 4-\frac{4}{3}
= \frac{8}{3}.
\end{align*}
}

\vfill 
\pagebreak

\subsection{Center of Mass ``Understand''}

The \textbf{moment about the y-axis} $M_y$ measures the tendency of an object or system to rotate about the y-axis. The \textbf{moment about the x-axis} $M_x$ measures the tendency of an object or system to rotate about the x-axis.


\begin{align*}
    \overline{x}&=\frac{1}{m}\int\int_D \, x \, \rho(x,y) \, dA\\
    \overline{y}&=\frac{1}{m}\int\int_D \, y \, \rho(x,y) \, dA
\end{align*}


\subsection{Example ``Do''}

Find the center of mass of a triangular lamina with vertices $(0,0),(1,0)$ and $(0,2)$ and density function $\rho(x,y)=1+3x+y$.

\sol{%
\textbf{Solution.}
From the previous example, $m=\dfrac{8}{3}$ and the inner integral of $\rho$ was $4-4x^2$ over $0\leq y\leq 2-2x$.
\begin{align*}
M_y &= \iint_D x\rho\,dA
= \int_0^1 x(4-4x^2)\,dx
= \int_0^1(4x-4x^3)\,dx
= \bigl[2x^2-x^4\bigr]_0^1
= 1,
\end{align*}
so $\displaystyle\overline{x}=\frac{M_y}{m}=\frac{1}{8/3}=\frac{3}{8}$.
\begin{align*}
M_x &= \iint_D y\rho\,dA
= \int_0^1\int_0^{2-2x} y(1+3x+y)\,dy\,dx
= \int_0^1\int_0^{h}(y+3xy+y^2)\,dy\,dx,
\end{align*}
where $h=2-2x$.
\begin{align*}
\Bigl[\frac{y^2}{2}+\frac{3x y^2}{2}+\frac{y^3}{3}\Bigr]_0^{h}
&= \frac{h^2}{2}(1+3x)+\frac{h^3}{3}
= 2(1-x)^2(1+3x)+\frac{8}{3}(1-x)^3.
\end{align*}
Expanding and integrating:
\begin{align*}
M_x &= \int_0^1\Bigl(\frac{14}{3}-6x-2x^2+\frac{10}{3}x^3\Bigr)\,dx
= \Bigl[\frac{14}{3}x-3x^2-\frac{2}{3}x^3+\frac{5}{6}x^4\Bigr]_0^1
= \frac{11}{6}.
\end{align*}
Thus $\displaystyle\overline{y}=\frac{M_x}{m}=\frac{11/6}{8/3}=\frac{11}{16}$.\\[0.4em]
Center of mass: $\displaystyle\Bigl(\frac{3}{8},\frac{11}{16}\Bigr)$.
}

\vfill 

\pagebreak

\section{Practice}

\subsection{}

Find the mass and $\overline{y}$ of a triangular lamina with vertices $(0,1),(1,1),$ and $(1,0)$ with density function $\rho(x,y)=xy$.

\sol{%
\textbf{Solution.}
The triangle has hypotenuse from $(0,1)$ to $(1,0)$, i.e.\ $y=1-x$. The region is $0\leq x\leq 1$, $1-x\leq y\leq 1$.
\begin{align*}
m &= \int_0^1\int_{1-x}^{1} xy\,dy\,dx
= \int_0^1 x\cdot\frac{1}{2}\bigl[y^2\bigr]_{1-x}^{1}\,dx
= \frac{1}{2}\int_0^1 x\bigl(1-(1-x)^2\bigr)\,dx \\
&= \frac{1}{2}\int_0^1 x\bigl(2x-x^2\bigr)\,dx
= \frac{1}{2}\int_0^1(2x^2-x^3)\,dx
= \frac{1}{2}\Bigl[\frac{2}{3}x^3-\frac{1}{4}x^4\Bigr]_0^1
= \frac{5}{24}.
\end{align*}
\begin{align*}
M_x &= \iint_D y\cdot\rho\,dA
= \int_0^1\int_{1-x}^{1} x y^2\,dy\,dx
= \int_0^1 x\cdot\frac{1}{3}\bigl[y^3\bigr]_{1-x}^{1}\,dx
= \frac{1}{3}\int_0^1 x\bigl(1-(1-x)^3\bigr)\,dx \\
&= \frac{1}{3}\int_0^1 x(3x-3x^2+x^3)\,dx
= \frac{1}{3}\int_0^1(3x^2-3x^3+x^4)\,dx
= \frac{1}{3}\Bigl(1-\frac{3}{4}+\frac{1}{5}\Bigr)
= \frac{3}{20}.
\end{align*}
\begin{align*}
\overline{y} &= \frac{M_x}{m} = \frac{3/20}{5/24} = \frac{3}{20}\cdot\frac{24}{5} = \frac{18}{25}.
\end{align*}
So $m=\dfrac{5}{24}$ and $\overline{y}=\dfrac{18}{25}$.
}

\vfill

\subsection{}%15.4 number 3

Find the total mass and center of mass of the lamina that occupies the region $D$ and has the given density function $\rho$.
\begin{align*}
    D=\{ (x,y) \mid 1 \leq x \leq 3, 1 \leq y \leq 4\}; \rho(x,y)=ky^2
\end{align*}

\sol{%
\textbf{Solution.}
\begin{align*}
m &= \int_1^3\int_1^4 ky^2\,dy\,dx
= k\int_1^3\Bigl[\frac{y^3}{3}\Bigr]_1^4\,dx
= k\int_1^3\frac{63}{3}\,dx
= 21k\cdot 2
= 42k.
\end{align*}
\begin{align*}
M_y &= \iint_D x\rho\,dA
= \int_1^3\int_1^4 xky^2\,dy\,dx
= k\int_1^3 x\cdot 21\,dx
= 21k\Bigl[\frac{x^2}{2}\Bigr]_1^3
= 21k\cdot 4
= 84k,
\end{align*}
so $\displaystyle\overline{x}=\frac{84k}{42k}=2$.
\begin{align*}
M_x &= \iint_D y\rho\,dA
= \int_1^3\int_1^4 ky^3\,dy\,dx
= k\int_1^3\Bigl[\frac{y^4}{4}\Bigr]_1^4\,dx
= k\int_1^3\frac{255}{4}\,dx
= k\cdot\frac{255}{4}\cdot 2
= \frac{255k}{2},
\end{align*}
so $\displaystyle\overline{y}=\frac{255k/2}{42k}=\frac{255}{84}=\frac{85}{28}$.\\[0.4em]
Mass $m=42k$; center of mass $\displaystyle\Bigl(2,\dfrac{85}{28}\Bigr)$.
}

\vfill

\pagebreak
\section{Harder Practice}

    \subsection{} Find the center of mass of a semicircular plate with of radius $r$. The plate has constant density. %See page 563. \\
    
\sol{%
\textbf{Solution.}
Take the upper half-disk $x^2+y^2\leq r^2$, $y\geq 0$, with constant density $\rho$. By symmetry, $\overline{x}=0$.\\[0.4em]
Mass: $m=\rho\cdot\dfrac{\pi r^2}{2}$.\\
In polar coordinates,
\begin{align*}
M_x &= \iint_D y\rho\,dA
= \rho\int_0^{\pi}\int_0^{r}(s\sin\theta)\,s\,ds\,d\theta
= \rho\int_0^{\pi}\sin\theta\,d\theta\int_0^{r} s^2\,ds
= \rho\cdot 2\cdot\frac{r^3}{3}
= \frac{2\rho r^3}{3}.
\end{align*}
\begin{align*}
\overline{y} &= \frac{M_x}{m}
= \frac{2\rho r^3/3}{\rho\pi r^2/2}
= \frac{2r^3}{3}\cdot\frac{2}{\pi r^2}
= \frac{4r}{3\pi}.
\end{align*}
Center of mass: $\displaystyle\Bigl(0,\dfrac{4r}{3\pi}\Bigr)$.
}

\vfill

    \subsection{} Find the center of mass of a semicircular plate whose density at any point is proportional to the distance from the center of the circle. 

\sol{%
\textbf{Solution.}
Again take the upper half-disk of radius $r$. Density $\rho(x,y)=k\sqrt{x^2+y^2}=ks$ in polar. By symmetry, $\overline{x}=0$.
\begin{align*}
m &= \int_0^{\pi}\int_0^{r}(ks)\,s\,ds\,d\theta
= k\int_0^{\pi}d\theta\int_0^{r} s^2\,ds
= k\pi\cdot\frac{r^3}{3}
= \frac{k\pi r^3}{3}.
\end{align*}
\begin{align*}
M_x &= \iint_D y\rho\,dA
= \int_0^{\pi}\int_0^{r}(s\sin\theta)(ks)\,s\,ds\,d\theta
= k\int_0^{\pi}\sin\theta\,d\theta\int_0^{r} s^3\,ds
= k\cdot 2\cdot\frac{r^4}{4}
= \frac{k r^4}{2}.
\end{align*}
\begin{align*}
\overline{y} &= \frac{M_x}{m}
= \frac{kr^4/2}{k\pi r^3/3}
= \frac{r^4}{2}\cdot\frac{3}{\pi r^3}
= \frac{3r}{2\pi}.
\end{align*}
Center of mass: $\displaystyle\Bigl(0,\dfrac{3r}{2\pi}\Bigr)$.
}


\vfill


\vspace{8mm}

\end{document}
